%==============================================================
% ICCIT 2026 submission -- IEEE conference format
% FULL VERSION: every section from the original draft is restored.
% This WILL exceed the 6-page ICCIT limit. Trim later; see the
% "TRIM CANDIDATES" list at the bottom of this file.
%
% Format requirements already applied: A4 paper, ANONYMOUS author block,
% IEEE numeric references, no ACM leftovers.
%
% Build:  pdflatex iccit2026_full && pdflatex iccit2026_full
%
% >>> BEFORE SUBMISSION, resolve every "TODO" comment below. <<<
%==============================================================
\documentclass[conference,a4paper]{IEEEtran}
\IEEEoverridecommandlockouts

\usepackage{cite}
\usepackage{amsmath,amssymb,amsfonts}
\usepackage{graphicx}
\usepackage{booktabs}
\usepackage{textcomp}
\usepackage{url}

\graphicspath{{figures/}}

%--------------------------------------------------------------
% TWO SWITCHES. Both are set for a DRAFT right now.
%
%   \blindsubmissiontrue  -> anonymous author block, Acknowledgment hidden
%   \blindsubmissionfalse -> real authors + Acknowledgment printed
%
%   \draftnotestrue  -> internal "verify before submission" notes PRINT
%   \draftnotesfalse -> those notes are hidden
%
% >>> SET \draftnotesfalse BEFORE YOU SUBMIT. Those notes were visible in
% >>> the original PDF's reference list and read as an unfinished paper.
%--------------------------------------------------------------
\newif\ifblindsubmission
\blindsubmissiontrue      % ICCIT 2026 requires anonymous submission

\newif\ifdraftnotes
\draftnotestrue           % <-- flip to \draftnotesfalse before submitting

% System name used throughout the original draft.
\newcommand{\sysname}{SignalCheck}

% TODO(verify): confirm the exact hosted model identifier actually used in
% the experiments. The original draft said "gemini-3.5-flash", which is not a
% released model name. Fix this macro once, it propagates everywhere.
\newcommand{\llmProprietary}{gemini-2.5-flash}
\newcommand{\llmOpen}{gpt-oss-120b}

\begin{document}

%--------------------------------------------------------------
% TITLE -- two options. The first is active.
%
% (A) ACTIVE: revised title. Drops the "BLIP-CLIP Feature Fusion" claim that
%     Table I does not support, and scopes the verification stage as a
%     prototype. This is the reviewer-4 / reviewer-5 fix.
%
% (B) ORIGINAL title from the submitted draft, kept verbatim below. Swap the
%     comment markers if you want it back -- but note a reviewer will compare
%     "Feature Fusion" in the title against CLIP-only matching the fused model
%     in Table I.
%--------------------------------------------------------------
\title{\sysname: Multimodal Bangla Fake News Image Detection --
Modality Ablation, Distribution-Shift Evaluation,\\
and an Evidence-Grounded Verification Prototype}

% \title{Multimodal Bangla Fake News Detection Through BLIP--CLIP
% Feature Fusion, XGBoost Classification, and Retrieval-Augmented\\
% Dual-LLM Verification}

%--------------------------------------------------------------
% AUTHOR BLOCK -- controlled by \blindsubmission (see preamble).
%--------------------------------------------------------------
\ifblindsubmission
  \author{\IEEEauthorblockN{Anonymous Submission}
  \IEEEauthorblockA{Paper ID: XXXX}}
\else
  \author{\IEEEauthorblockN{1\textsuperscript{st} Given Name Surname}
  \IEEEauthorblockA{\textit{dept. name of organization} \\
  \textit{name of organization}\\ City, Country \\ email or ORCID}
  \and
  \IEEEauthorblockN{2\textsuperscript{nd} Given Name Surname}
  \IEEEauthorblockA{\textit{dept. name of organization} \\
  \textit{name of organization}\\ City, Country \\ email or ORCID}
  \and
  \IEEEauthorblockN{3\textsuperscript{rd} Given Name Surname}
  \IEEEauthorblockA{\textit{dept. name of organization} \\
  \textit{name of organization}\\ City, Country \\ email or ORCID}}
\fi

\maketitle

\begin{abstract}
Fake news detection for low-resource languages such as Bangla is held back by
two things: the scarcity of large-scale, image-centric training data, and
classifiers that report only same-distribution accuracy without any external
evidence-verification stage. This paper presents \sysname, a two-stage
multimodal pipeline for detecting manipulated news imagery circulating on
Bangladeshi social media. The first stage fuses a BLIP-generated image caption, encoded
via multilingual SBERT and refined through Cross-Image Semantic Fusion (CISF),
with a CLIP visual embedding, and classifies the combined feature vector with
a gradient-boosted decision tree. A systematic per-modality ablation on the
BanFakeMM benchmark shows that CLIP visual features dominate classification
accuracy, reaching 95.58 percent alone, and that no text-inclusive
configuration exceeds this beyond fold-to-fold variation. OCR-based text
extraction, found to perform near chance level in isolation at 56.80 percent,
is removed from the final configuration. Evaluated on a disjoint AdversarialGen
manipulation style, zero-shot accuracy falls to 86.0 percent with fake-recall
at 73.0 percent, revealing a substantial distribution-shift gap; threshold
adjustment and supervised training-data diversification raise accuracy to
91.5 percent, which we report explicitly as domain adaptation rather than
unseen-style generalization. The second stage is a retrieval-augmented,
dual-independent-LLM verification layer grounded in a Bangla-language
fact-check evidence corpus, reaching a supported, contradicted or insufficient
verdict either from claim text or directly from CLIP-based image-to-image
evidence matching. We report the classifier's ablation and distribution-shift
results quantitatively, characterize the verification layer's behaviour
through qualitative case studies, present it explicitly as a prototype, and
discuss the honest limitations of both components.
\end{abstract}

%--------------------------------------------------------------
% CCS Concepts -- this is an ACM classification block and has NO IEEE
% equivalent; IEEE proceedings use Index Terms only. It is reproduced here
% because it was present in the original draft. Most IEEE reviewers will read
% it as a leftover from another template, so consider deleting it.
%--------------------------------------------------------------
\begin{IEEEkeywords}
\textbf{CCS Concepts:} Computing methodologies $\rightarrow$ Machine learning
approaches; Natural language processing; Information systems $\rightarrow$
Multimedia and multimodal retrieval.
\end{IEEEkeywords}

\begin{IEEEkeywords}
Bangla fake news detection, multimodal misinformation, BLIP, CLIP, XGBoost,
retrieval-augmented generation, dual-LLM verification, distribution shift,
domain adaptation, out-of-distribution evaluation, low-resource language
processing
\end{IEEEkeywords}

%==============================================================
\section{Introduction}
%==============================================================

Bangladesh's social media ecosystem---dominated by Facebook, YouTube, and
messaging-app forwards---has become a primary vector for manipulated and
fabricated news imagery, ranging from digitally altered photocards and forged
documents to fully AI-generated video frames misattributed to public figures.
Unlike text-only misinformation, this class of fake news exploits the combined
visual and textual channel: a viral image is rarely judged on its pixels
alone, but on the claim, caption, or context attached to it as it spreads.

This problem is compounded by linguistic and cultural specificity. Bangla,
unlike English, is a low-resource language for NLP, so most fake-news
detection methods---sentiment, ensemble, graph, and comment-fusion models
alike---have been built and tested only on English datasets from Western
platforms such as LIAR, WELFake and FakeNewsNet. These systems presuppose
signal sources such as dense comment threads or public retweet graphs that are
frequently sparse, unavailable, or structured differently in the Bangladeshi
context. This research identified several core problems in the current state
of Bangla multimodal fake-news detection.

\subsection{Research Problems}

\begin{itemize}
\item \textbf{Training-data scarcity.} No large-scale, organically collected,
labeled dataset of manipulated Bangladeshi news images circulating in
real-world settings exists; training must rely on synthetically constructed
benchmarks such as BanFakeMM.

\item \textbf{Claim-caption divergence.} An uploaded image's automatically
generated caption describes what is visually present in it, not the claim
being asserted about it, so a verdict cannot be reliably reached from the
image alone.

\item \textbf{Evidence-source scarcity.} Bangla-language, regionally relevant
fact-checking corpora are far sparser than the Wikipedia-scale corpora used in
general-purpose retrieval-augmented systems.

\item \textbf{No integrated pipeline.} No prior Bangla-focused system combines
image-level manipulation detection with an external evidence-verification
stage in a single pipeline.

\item \textbf{Same-distribution overstatement.} Same-distribution accuracy is
routinely reported as if it were real-world reliability, without a disjoint
evaluation to check whether a classifier has merely learned one narrow
manipulation style.
\end{itemize}

\subsection{Approach and Contributions}

This research addresses these problems as follows.

\begin{itemize}
\item \textbf{Per-modality ablation.} Deciding which modalities---caption
text, OCR text, CLIP visual features, CISF fusion---actually contribute,
rather than including every available signal by default, and empirically
removing OCR once ablation showed it performs near chance level in isolation
on this dataset. Results are reported with precision, recall and F1 alongside
accuracy.

\item \textbf{Evaluation under distribution shift.} Using a held-out,
non-overlapping BanFakeMM manipulation-style split to quantify the model's
generalization gap, then narrowing it through threshold calibration and
targeted training-data diversification. We separate explicitly what is
recovered by thresholding from what is recovered by \emph{supervised
adaptation} on target-style data, and label the latter as domain adaptation
rather than unseen-style generalization.

\item \textbf{Retrieval-augmented, dual-independent-LLM verification.}
Grounding claim assessment in an actual Bangla-language fact-check evidence
corpus rather than trusting the classifier's risk score, or a single LLM's
unverified judgment, as a final answer. This stage is presented as a
\emph{prototype} pending a labelled quantitative evaluation.

\item \textbf{CLIP-based image-to-image evidence matching.} Comparing the
uploaded image's own pixels against previously fact-checked images, bypassing
unreliable caption-derived claims whenever the exact image has already been
independently verified elsewhere.

\item \textbf{Explicit support for user-supplied claim text.} Letting a user
provide the actual circulating caption or claim, since an automatic image
caption alone cannot state what is being claimed about an image.
\end{itemize}

Unlike previous research, this study evaluates its classifier not only on a
same-distribution held-out split but on a disjoint manipulation-style split
never seen during training, and pairs that classifier with an
evidence-verification stage that neither prior Bangla multimodal detectors nor
most general-purpose retrieval-augmented fact-checking systems combine with
image-level manipulation classification in a single pipeline.

While BLIP/CLIP-based multimodal classification, dual-LLM verification
architectures, and retrieval-augmented fact-checking have each been explored
independently in prior work, there is, to our knowledge, no existing study
that combines image-level manipulation detection, a genuinely disjoint
distribution-shift evaluation, and a retrieval-augmented dual-LLM
verification stage---including direct image-to-image evidence matching---within
a single pipeline for a low-resource language such as Bangla. This is the gap
this research addresses. We deliberately do not claim architectural novelty
for the fusion or verification paradigms themselves; the regional adaptation
and the evaluation protocol are where we locate this work's contribution.

%==============================================================
\section{Literature Review}
%==============================================================

Research on Bangla-language fake news detection remains comparatively limited
relative to English. BanFakeNews \cite{banfakenews} established one of the
earliest large-scale benchmarks for this task: an annotated corpus of roughly
50{,}000 Bangla news items labeled authentic or fabricated, introduced
specifically because most fake-news detection methods have been developed for
English while low-resource languages such as Bangla remain comparatively
under-served. Subsequent work extended this line to multimodal settings.
MultiBanFakeDetect \cite{multibanfake} combines textual and visual fusion
techniques, reporting a hybrid DenseNet-169/mBERT fusion architecture to
detect Bangla fake news in under-resourced contexts, and is, to our knowledge,
the closest prior system to the present work in both language and multimodal
scope. Beyond Bangla specifically, the MMCFND framework and its accompanying
MMIFND dataset \cite{mmcfnd} extend multimodal, caption-aware fake-news
detection to seven low-resource Indic languages, including Bengali, with
roughly 28{,}000 labeled instances, and BanFakeNews-2.0 \cite{banfakenews2}
introduced a 60{,}000-instance manually curated Bangla text corpus
(47{,}000 authentic, 13{,}000 fake).

None of these systems, however, is built around a defined image-manipulation
task with adversarially constructed positive examples---MultiBanFakeDetect
labels whole news articles by their overall claim, and BanFakeNews-2.0 is
text-only---nor does any of them incorporate an external evidence-retrieval or
claim-verification component. Outside the Bangla-specific literature, EANN
\cite{eann}, MVAE \cite{mvae}, and SpotFake \cite{spotfake} each established
that visual and textual signals are complementary for fake-news
classification, but each was developed and validated on English-language,
Western-platform datasets, and none pairs classification with a
retrieval-augmented evidence-verification stage.

A more recent thread of work pairs multimodal classification with
retrieval-augmented, multi-model evidence verification directly. DEFAME
\cite{defame} orchestrates multiple specialized LLM-based experts against
retrieved multimodal evidence; RAGAR \cite{ragar} applies retrieval-augmented
reasoning to political fact-checking with multimodal large language models;
RAMA \cite{rama} proposes a retrieval-augmented multi-agent framework for
multimodal misinformation detection; and CaLM \cite{calm} contrasts a large
and a small language model to verify grounded generation. None of these is
Bangla-specific or Bangladesh-corpus-grounded.

This research identified four critical limitations in the surveyed literature.

\begin{itemize}
\item \textbf{Modality-scope limitation.} Bangla-focused systems are trained
and evaluated on curated news-article datasets rather than the screenshot- and
photocard-style imagery commonly observed circulating on Bangladeshi social
media.

\item \textbf{Verification-scope limitation.} No Bangla-focused system, to our
knowledge, incorporates an external evidence-retrieval or claim-verification
stage alongside image-level classification.

\item \textbf{Generalization-reporting limitation.} General-purpose
retrieval-augmented verification systems such as DEFAME, RAGAR and RAMA are
not Bangla-specific, and Bangla multimodal detectors do not report a disjoint
distribution-shift evaluation.

\item \textbf{Novelty-attribution limitation.} Prior work sometimes frames
dual-LLM or retrieval-augmented verification as architecturally novel, when
the paradigm itself is already explored in DEFAME, RAGAR, RAMA and CaLM---the
regional adaptation, not the architecture, is where genuine novelty should be
claimed.
\end{itemize}

This research addresses all four limitations simultaneously: it targets
image-level manipulation detection on screenshot/photocard-style imagery via
BanFakeMM, evaluates that classifier under a genuine distribution-shift
protocol, adds a Bangla-grounded retrieval-augmented dual-LLM verification
stage with direct image-to-image evidence matching, and explicitly scopes its
novelty claim to the regional adaptation and integration rather than the
underlying verification paradigm.

%==============================================================
\section{Methodology}
%==============================================================

\subsection{System Architecture}

\sysname{} consists of two cooperating stages. The first is a local multimodal
classifier that assigns a fake-probability score to an uploaded image using
only its visual and captioned content. The second is an evidence-grounded
verification layer that retrieves independently sourced fact-check
evidence---matched either by the image's own pixels or by its associated claim
text---and produces an evidence-based verdict using two independently
configured large language models. The two stages are complementary: the
classifier provides a fast, dataset-trained risk signal even when no external
evidence exists, while the verification layer provides claim-level,
provenance-backed reasoning whenever relevant evidence can be found.
Fig.~\ref{fig:arch} illustrates the complete pipeline, including the
retrieval-augmented generation boundary within the verification stage.

\begin{figure}[t]
\centering
\includegraphics[width=\columnwidth]{fig_architecture.png}
\caption{\sysname{} system architecture---local multimodal classifier (top) and
retrieval-augmented dual-LLM verification layer (bottom).}
\label{fig:arch}
\end{figure}

\subsection{Dataset and Preprocessing}

The classifier is trained and evaluated on the AdversarialEdit split of the
public BanFakeMM benchmark \cite{banfakemm}: 2{,}000 images, 1{,}000 authentic
and 1{,}000 manipulated, giving a class-balanced base dataset. BanFakeMM
constructs its ``fake'' images through a defined editing or generation
pipeline rather than through subjective human annotation of ambiguous
real-world content, so ground-truth labels in this dataset are deterministic
by construction. To assess generalization beyond this training distribution, a
disjoint evaluation set is drawn from BanFakeMM's AdversarialGen split, a
fully AI-generated fake-image collection with no image-pool overlap with the
training data.

\subsubsection{Preprocessing steps}

\begin{itemize}
\item \textbf{Captioning and OCR.} A BLIP model generates a natural-language
caption describing each image's visual content. OCR-based text extraction was
evaluated during an earlier ablation phase and found to perform near the
chance level in isolation on this dataset (56.80 percent accuracy); since the
final CLIP-based configuration achieved comparable or better accuracy without
it, OCR was excluded from the final preprocessing pipeline.

\item \textbf{Text processing and embedding.} The caption is tokenized and
stemmed, with separate handling for English and Bangla script, to produce a
processed-text representation used for TF-IDF features, and separately encoded
into a dense semantic embedding via a multilingual SBERT model.

\item \textbf{Visual embedding.} A CLIP vision encoder independently extracts
a 512-dimensional visual embedding directly from the image pixels.

\item \textbf{Shifted evaluation set construction.} 200 AdversarialGen images
(100 authentic, 100 manipulated) are reserved for the shifted evaluation and
are never used in training; a separate, non-overlapping sample of 300 images
is used only for the training-data diversification experiment in
Section~\ref{sec:shift}.

\item \textbf{Evidence corpus construction.} The verification layer's evidence
corpus is built independently of the classifier's training data, from a single
Bangla-language fact-checking publisher (Rumor Scanner), via a sitemap-based
and keyword-targeted scraping process. Of the 1{,}521 documents indexed,
1{,}465 (96 percent) include the article's full body text; the remaining 56
(4 percent) fall back to a headline-level summary.
\end{itemize}

% TODO(reviewer-2, BLOCKING): disjointness is currently established by
% filename comparison only. Before submission, re-verify with a perceptual
% hash (pHash/dHash) AND a CLIP-embedding near-duplicate sweep across all
% partitions, report the number of duplicate clusters removed, and rebuild
% the cross-validation folds with GroupKFold grouped by duplicate cluster and
% source article. Then delete the disclosure paragraph below and state the
% verified result instead.
\emph{Disclosure: duplicate verification.} Partition disjointness is currently
established by filename comparison rather than by perceptual-hash or
embedding-based near-duplicate detection, and cross-validation folds are
assigned at the image level by stratified randomization rather than grouped by
source article or near-duplicate cluster. Residual near-duplicate leakage is
therefore an unresolved risk rather than a ruled-out one, and the
same-distribution figures below should be read with that caveat.

\subsection{Multimodal Feature Fusion and Classification}

For each input image, three complementary representations are extracted. The
BLIP caption is encoded into a semantic embedding using SBERT, then refined
through Cross-Image Semantic Fusion (CISF): the embedding is combined with the
weighted average of its $k=5$ nearest semantic neighbours within the training
corpus, using a fusion weight of $\alpha=0.65$ favouring the original
embedding, and the fused vector is re-normalized and reduced to 32 principal
components via PCA. In parallel, the CLIP visual embedding (512-dimensional)
is separately reduced to 64 principal components via its own PCA transform,
fit only on the training fold in the same way as the caption pathway's
reduction---this keeps the two embedding blocks on a comparable scale relative
to the lower-dimensional TF-IDF and numeric features.

The neighbour count ($k=5$), fusion weight ($\alpha=0.65$), and target
dimensionality (32 components) were fixed design choices for this study rather
than selected through a dedicated validation search; a systematic sensitivity
study over these hyperparameters, and comparison against adaptive fusion
mechanisms such as MMFusionNet \cite{mmfusionnet} or CLAAF \cite{claaf}, is
left for future work. To prevent data leakage, each training embedding's own
nearest-neighbour search excludes itself when the fusion reference set is
constructed, and at inference time an unseen image's embedding is compared
only against this fixed, pre-established training-set reference.

The final feature vector concatenates four blocks, each fit only on the
training fold: the 32-component CISF-fused caption embedding, the 64-component
CLIP embedding, a TF-IDF representation over the processed caption text
(unigrams and bigrams, vocabulary capped at 5{,}000 features, minimum document
frequency 3), and a small set of auxiliary numeric features (token counts,
neighbour-similarity statistics). This combined representation is passed to a
gradient-boosted decision tree classifier (XGBoost). Ablation, reported in
Section~\ref{sec:ablation}, shows the CLIP block to be the dominant
contributor to classification accuracy, with the fused caption embedding
contributing complementary rather than additive signal.

\subsection{Training Configuration}
\label{sec:training}

XGBoost hyperparameters were selected via randomized search (scikit-learn's
\texttt{RandomizedSearchCV}, 8 sampled configurations evaluated under 3-fold
cross-validation, scored by fake-class F1) rather than manual choice---a
comparatively small search budget chosen for computational cost, with a wider
budget left for future tuning. The resulting configuration is used throughout
this paper's reported results:

\begin{itemize}
\item \texttt{max\_depth}~3, \texttt{learning\_rate}~0.0742,
\texttt{n\_estimators}~443
\item \texttt{subsample}~0.866, \texttt{colsample\_bytree}~0.683
\item \texttt{min\_child\_weight}~2, \texttt{gamma}~0.143
\item \texttt{reg\_alpha}~$1.01\times10^{-5}$, \texttt{reg\_lambda}~4.911
\end{itemize}

%--------------------------------------------------------------
% THRESHOLD PARAGRAPH -- two versions.
%
% (A) ACTIVE below: describes what was ACTUALLY done. The threshold was swept
%     on the evaluation set, so this paragraph states that plainly and defers
%     to the disclosure in Sec. IV-C. Honest, but the protocol is still leaky.
%
% (B) COMMENTED OUT further down: the CORRECTED protocol. Do NOT activate it
%     until you have actually run it -- it describes a 300/150/200 partition
%     that does not exist yet. Activating it early would claim a methodology
%     that was not followed, which is worse than the disclosed leak.
%
%     Activation checklist (all must be true):
%       [ ] pHash + CLIP-embedding dedup run; duplicate clusters identified
%       [ ] AdversarialGen re-partitioned 300 / 150 / 200 at cluster level
%       [ ] selection criterion fixed BEFORE sweeping (argmax fake-class F1)
%       [ ] tau swept on the 150-image validation partition only
%       [ ] 200-image sealed partition scored exactly once
%       [ ] [X] below replaced with the real value
%       [ ] Table II daggered rows replaced; Sec. IV-C disclosure deleted
%--------------------------------------------------------------
The classifier outputs a fake-probability score in $[0,1]$. The operating
threshold used throughout this paper is 0.30 rather than the conventional
0.50, because that value substantially improves fake-recall at a modest cost
to real-image precision on the shifted evaluation set. This threshold was
selected by sweeping that same evaluation set rather than a disjoint
validation partition; the resulting selection effect, and which reported
figures it does and does not affect, are disclosed in
Section~\ref{sec:shift}.

% ---------- (B) CORRECTED PROTOCOL -- activate only after re-running ----------
% The classifier outputs a fake-probability score in $[0,1]$. Rather than
% adopting the conventional 0.50 cut-off, the operating threshold is selected
% on held-out validation data and never on the evaluation set. Two thresholds
% are reported. For same-distribution evaluation, $\tau_{\mathrm{in}}$ is
% chosen on the inner validation folds of the base-split cross-validation. For
% the distribution-shift evaluation, the AdversarialGen split is partitioned at
% the near-duplicate-cluster level into three disjoint subsets: 300 images for
% adaptation training, 150 for threshold-selection validation, and 200 sealed
% for final testing. $\tau_{\mathrm{adapt}} = [X]$ is the value maximizing
% fake-class F1 on the 150-image validation partition; this criterion was fixed
% prior to the sweep. The 200-image test partition is evaluated exactly once,
% after all threshold and training decisions are frozen, and is used for no
% tuning of any kind.
% -----------------------------------------------------------------------------

Per-modality ablation results (Section~\ref{sec:ablation}) come from repeated
stratified 5-fold cross-validation (2 repeats per configuration, 10 folds
total), with identical fold splits and fixed hyperparameters held constant
across every evaluated feature configuration. Folds are constructed at the
image level via stratified random assignment; they are not explicitly grouped
by source article or near-duplicate image, which is noted as a limitation in
Section~\ref{sec:limitations}. All dimensionality-reduction transforms,
vocabularies and neighbour reference sets are fit inside the training fold
only.

% TODO(reviewer-6): fill in the exact values actually used, then delete this
% comment. These are required for the reproducibility paragraph to be true.
\emph{Reproducibility.} All runs use a fixed random seed
(\texttt{seed}~=~[SEED]) for fold assignment, the randomized search and the
classifier. Software versions, exact BLIP/CLIP/SBERT checkpoint identifiers,
hardware and end-to-end runtime are given in the supplementary material, and
code and feature files will be released on acceptance.

\subsection{Retrieval-Augmented Dual-LLM Verification}
\label{sec:verify}

The claim to be verified is determined by priority: a user-supplied claim
text, if provided, is used directly; otherwise the BLIP caption is used as a
fallback. In parallel, the uploaded image's own CLIP embedding is compared
against the embeddings of images from previously fact-checked articles; a
similarity threshold of 0.90 is used to treat a match as strong evidence that
this exact image has already been independently verified---unrelated image
pairs typically scored 0.55--0.60 on CLIP cosine similarity during
development, while genuine re-shared copies scored 0.90 or above, though this
threshold reflects an approximate empirical separation rather than a formally
calibrated precision/recall curve.

When a direct match is found, the system reports the matched article and its
similarity score as evidence rather than asserting a final truth verdict on
its own, since an image match establishes only that the image has previously
circulated and been fact-checked, not that the claim currently attached to it
is the same claim the matched article addressed. When no direct match exists,
the resolved claim text queries the local evidence store via SBERT-based
semantic similarity search (top-5 matches, minimum similarity 0.30), and,
separately, live web-search results are retrieved. Both are passed to two
independently configured large language models---\texttt{\llmProprietary} and
the open-weight \texttt{\llmOpen}, served through Groq's inference API at
temperature 0---each given an identical prompt containing the claim and the
retrieved evidence, instructed to return a JSON verdict (\emph{supported},
\emph{contradicted}, or \emph{insufficient}) with a confidence score, a
rationale, and cited source URLs.

Each model is queried independently, with no visibility into the other's
output. A consensus verdict is returned only when both models complete
successfully and agree; if either model errors, is not configured, or the two
models disagree, the system withholds a consensus verdict and flags the case
as needing human review, rather than defaulting to whichever single response
is available.

% TODO(reviewer-5, BLOCKING if the prototype framing is dropped): to promote
% this from prototype to evaluated component, build a labelled claim set
% (target 150-200 claims with adjudicated supported/contradicted/insufficient
% labels), then report: retrieval recall@1 and recall@5, verdict accuracy and
% macro-F1, abstention rate, image-match precision at the 0.90 threshold, a
% table of 2-3 worked examples, and a single-model (each model alone) versus
% dual-LLM consensus comparison. RAGProbe [17] and the RAG metric survey
% [13] are cited below as concrete starting points. Until those exist, the
% scoping paragraph below must stay.
\emph{Scope.} Both the 0.90 similarity threshold and the consensus rule are
practical operating points set from informal development observation, not from
a calibrated precision/recall study. No labelled claim set, retrieval metric,
verdict-accuracy measurement, or single-versus-dual-LLM comparison is
reported in this paper, and no sustained deployment trial was conducted.
Accordingly this stage is presented as a working prototype that demonstrates
feasibility and defines an evaluation target, not as a validated verification
system, and none of the quantitative claims in this paper depend on it.

\subsection{Evaluation Metrics}

Classification results are reported as accuracy, precision, recall and F1 for
the manipulated class, with real-class recall additionally reported for the
distribution-shift evaluation, since accuracy alone can mask asymmetric
degradation between classes under shift. Fake-recall (fakes correctly
identified over total actual fakes) is treated as the more safety-critical
metric, since a missed fake image can circulate unchecked whereas a false
alarm on a real image costs only a review step; this asymmetric cost is why
the deployed decision threshold favours fake-recall. Confusion matrices are
reported for the shifted evaluation to make this trade-off directly visible.

For the evidence-verification layer, no formal quantitative metric is defined
in this paper; verdicts are instead assessed through qualitative case studies
covering both evidence-available and evidence-unavailable scenarios, and a
complete quantitative evaluation---retrieval recall@$k$, verdict accuracy
against ground-truth-labeled claims---is identified as future work in
Section~\ref{sec:limitations}.

%==============================================================
\section{Experimental Results}
%==============================================================

\subsection{Classifier Ablation Performance}
\label{sec:ablation}

\begin{table}[t]
\caption{Per-Modality Ablation on the Base Split, Repeated Stratified 5-Fold
CV ($2$ Repeats, $10$ Folds). Accuracy Shown as Mean $\pm$ Standard Deviation
Across Folds. All Values in Percent.}
\label{tab:ablation}
\centering
\setlength{\tabcolsep}{3.5pt}
\begin{tabular}{lcccc}
\toprule
\textbf{Configuration} & \textbf{Acc.} & \textbf{Prec.$_{f}$} &
\textbf{Rec.$_{f}$} & \textbf{F1$_{f}$} \\
\midrule
OCR only              & 56.80 $\pm$ 1.78 & 63.44 & 32.55 & 42.94 \\
Caption only          & 74.40 $\pm$ 1.79 & 75.05 & 73.15 & 74.07 \\
Text only             & 76.55 $\pm$ 2.01 & 76.75 & 76.25 & 76.48 \\
Text + CISF           & 77.48 $\pm$ 1.37 & 77.43 & 77.65 & 77.50 \\
\midrule
CLIP only             & \textbf{95.58} $\pm$ 0.92 & 95.12 & 96.10 & 95.59 \\
Text + CLIP           & 95.43 $\pm$ 0.83 & 95.16 & 95.75 & 95.44 \\
CLIP + CISF           & 95.58 $\pm$ 1.01 & 95.25 & 95.95 & 95.59 \\
Caption + CISF + CLIP & 95.40 $\pm$ 0.97 & 95.10 & 95.75 & 95.41 \\
\bottomrule
\multicolumn{5}{l}{\footnotesize Subscript $f$ denotes the manipulated
(``fake'') class.}
\end{tabular}
\end{table}

\begin{figure}[t]
\centering
\includegraphics[width=\columnwidth]{fig1_ablation.png}
\caption{Per-modality ablation accuracy (repeated stratified 5-fold CV).}
\label{fig:ablation}
\end{figure}

Table~\ref{tab:ablation} and Fig.~\ref{fig:ablation} report accuracy for each
feature configuration under the repeated stratified cross-validation protocol.
OCR text used in isolation performs poorly, 56.80 $\pm$ 1.78 percentage
points, meaningfully close to the 50 percent chance baseline for this balanced
binary task, with fake-recall of only 32.55 percent, while CLIP visual
features alone already reach 95.58 $\pm$ 0.92 points. Combining caption text
with CLIP (95.43 $\pm$ 0.83 points) does not exceed CLIP alone, and adding
CISF on top of CLIP (95.58 $\pm$ 1.01 points) shows no measurable improvement.
These differences are small relative to each configuration's fold-to-fold
variability---the 0.18-point gap between CLIP-only and the final configuration
sits well within one fold's typical variation---so this should be read as
suggestive of visual-feature dominance on this specific dataset, not as a
statistically established or generalizable claim about caption text's value
for Bangla misinformation detection broadly.

All text-only configurations saturate below 78 percent, while any
configuration containing CLIP exceeds 95 percent. Visual features, not caption
text or its derived numeric properties, carry the discriminative signal for
this task.

% TODO(reviewer-4, BLOCKING): the fusion architecture is not currently
% justified by these numbers. Before submission do ONE of:
%   (a) demonstrate a measurable benefit -- run CLIP-only vs. fusion on the
%       shifted set of Table II, report a paired test (5x2cv F-test or
%       bootstrap 95% CI on the difference), and add calibration (ECE/Brier)
%       and AUROC; if fusion wins on robustness or calibration, say so and
%       keep it; OR
%   (b) simplify -- make CLIP + XGBoost the proposed model, retitle, and
%       demote CISF to a reported negative result.
% The paragraph below is honest but is NOT a substitute for either.
\emph{On the fusion design.} On same-distribution accuracy the caption pathway
earns its place only if it contributes robustness or calibration rather than
accuracy, which Table~\ref{tab:ablation} does not measure. Retaining a
multimodal architecture is not by itself a justification. We therefore report
the final configuration's parity with CLIP-only as an open question about the
fusion design rather than as a positive result, and treat CLIP-only as the
operative baseline any fusion must beat.

\subsection{Final Configuration Evaluation}

The selected final configuration---caption text fused via CISF, combined with
CLIP features, without OCR---reaches 95.40 $\pm$ 0.97 percentage points
same-distribution accuracy, with fake-class precision 95.10 percent, recall
95.75 percent and F1 95.41 percent. This is comparable to, though marginally
below, the CLIP-only ceiling, while retaining a multimodal architecture whose
justification remains open as noted above. These figures are cross-validated
means rather than a single held-out confusion matrix, since the base-split
evaluation protocol (Section~\ref{sec:training}) uses repeated $k$-fold
cross-validation rather than one fixed train/test split; per-fold true and
false-positive counts were not separately retained, so a single confusion
matrix for the base split is not reported here.

\subsection{Performance Under Distribution Shift}
\label{sec:shift}

\begin{table}[t]
\caption{Evaluation on the Disjoint AdversarialGen Manipulation Style
(200 Images, 100/100 Balanced). All Values in Percent.}
\label{tab:shift}
\centering
\setlength{\tabcolsep}{3pt}
\begin{tabular}{llccc}
\toprule
\textbf{Setting} & \textbf{$\tau$} & \textbf{Acc.} &
\textbf{Rec.$_{f}$} & \textbf{Rec.$_{r}$} \\
\midrule
\multicolumn{5}{l}{\emph{Zero-shot transfer (no target-style data in training)}} \\
Default threshold              & 0.50 & 86.0 & 73.0 & 99.0 \\
Adjusted threshold$^{\dagger}$ & 0.30 & 90.0 & 86.0 & 94.0 \\
\midrule
\multicolumn{5}{l}{\emph{Supervised domain adaptation ($+300$ AdversarialGen images)}} \\
Adapted model$^{\dagger}$      & 0.30 & 91.5 & 91.0 & 92.0 \\
\bottomrule
\multicolumn{5}{p{0.92\columnwidth}}{\footnotesize $^{\dagger}$Threshold
selected by sweeping this same evaluation set. These two rows are optimistic
and are not held-out estimates; see the disclosure below. Subscripts $f$ and
$r$ denote the manipulated and authentic classes.}
\end{tabular}
\end{table}

\begin{figure}[t]
\centering
\includegraphics[width=\columnwidth]{fig2_ood_progression.png}
\caption{Accuracy, fake-recall and real-recall across the three evaluation
settings of Table~\ref{tab:shift}.}
\label{fig:progression}
\end{figure}

Table~\ref{tab:shift} and Fig.~\ref{fig:progression} summarize the three-stage
evaluation on the 200-image, class-balanced AdversarialGen set. At the
original decision threshold, accuracy falls to 86.0 percent from the
95.40 percent same-distribution result, with fake-recall dropping sharply to
73.0 percent (27 of 100 held-out fakes missed)---the model, trained on a
narrow manipulation style, under-detects shifted fakes even though real-image
recall remains high at 99.0 percent. The 23-point fake-recall collapse, not
the 9.4-point accuracy drop, is the operationally meaningful number.

Sweeping the decision threshold, without retraining, raises accuracy to
90.0 percent at threshold 0.30 with fake-recall of 86.0 percent. Retraining
with an additional 300 diverse AdversarialGen images, evaluated at this
threshold with the 200-image evaluation set held out throughout, raises final
accuracy to 91.5 percent.

\emph{This second result is domain adaptation, not generalization to an unseen
style.} Once target-style images enter training, the evaluation no longer
measures zero-shot transfer, and the 91.5 percent figure must not be read as
evidence that the generalization gap has closed. Only the 86.0 percent row
measures unseen-style behaviour. The corresponding claim in an earlier draft
of this work---that diversification ``narrows the out-of-distribution gap to
91.5 percent''---conflates the two regimes and is withdrawn.

% TODO(reviewer-1, BLOCKING): remove the threshold-selection leakage.
% Re-partition AdversarialGen at the near-duplicate-cluster level into three
% disjoint subsets: ~300 adaptation-training / ~150 validation / ~200 sealed
% test. Fix the selection criterion in advance (argmax fake-class F1 on the
% validation partition), select tau there, then evaluate the sealed test
% partition EXACTLY ONCE. Replace the two daggered rows and delete the
% disclosure paragraph below. Expect the numbers to fall; that is the point.
\emph{Disclosure: threshold selection.} The threshold of 0.30 was selected by
sweeping the same 200-image set on which the daggered rows are reported.
Fig.~\ref{fig:sweep} shows that sweep. Those two rows are therefore an
optimistic upper bound on what threshold adjustment achieves, not held-out
estimates, and the true held-out values are expected to be lower. Only the
zero-shot default-threshold row (86.0 percent) is free of this selection
effect, and it is the only row we rely on for the paper's claims about
distribution shift. A protocol using a validation partition disjoint from a
sealed test partition is required to report the remaining rows as
measurements, and is the immediate next step for this work.

\begin{figure}[t]
\centering
\includegraphics[width=\columnwidth]{fig4_threshold_sweep.png}
\caption{Decision-threshold sweep on the shifted evaluation set. This sweep is
the source of the selection effect disclosed in Section~\ref{sec:shift}.}
\label{fig:sweep}
\end{figure}

\subsection{Confusion Matrix Analysis Under Shift}

\begin{figure}[t]
\centering
\includegraphics[width=\columnwidth]{fig3_confusion_matrices.png}
\caption{Confusion matrices before ($\tau=0.50$, original) and after
($\tau=0.30$, with diverse training) mitigation.}
\label{fig:cm}
\end{figure}

Fig.~\ref{fig:cm} gives the confusion matrices before and after mitigation. At
the original threshold, the classifier misclassifies only 1 of 100 real images
but misses 27 of 100 fakes---a strongly asymmetric error pattern consistent
with a decision boundary calibrated tightly around the training distribution's
specific manipulation signature. After threshold tuning and diverse-data
retraining, the matrix is substantially more balanced: 91 of 100 fakes
correctly caught (9 missed) and 92 of 100 real images correctly passed
(8 misclassified). This 91.5 percent result remains below the 95.40 percent
same-distribution ceiling, and the residual gap is read as an honest bound on
the classifier's demonstrated behaviour under shift, not as evidence that the
gap has fully closed.

\subsection{Baseline Comparison}

\begin{table}[t]
\caption{Baseline Comparison Under the Identical Cross-Validation Protocol.
All Values in Percent.}
\label{tab:baselines}
\centering
\setlength{\tabcolsep}{3.5pt}
\begin{tabular}{lcccc}
\toprule
\textbf{Model} & \textbf{Acc.} & \textbf{Prec.$_{f}$} &
\textbf{Rec.$_{f}$} & \textbf{F1$_{f}$} \\
\midrule
Linear SVM (TF-IDF)     & 76.30 $\pm$ 1.96 & 75.31 & 78.30 & 76.77 \\
XGBoost (TF-IDF + num.) & 78.95 $\pm$ 1.35 & 79.06 & 78.80 & 78.91 \\
\midrule
CLIP only               & 95.58 $\pm$ 0.92 & 95.12 & 96.10 & 95.59 \\
Proposed (final)        & 95.40 $\pm$ 0.97 & 95.10 & 95.75 & 95.41 \\
\bottomrule
\end{tabular}
\end{table}

\begin{figure}[t]
\centering
\includegraphics[width=\columnwidth]{fig5_baselines.png}
\caption{Multimodal configuration against text-only baselines under the
identical cross-validation protocol.}
\label{fig:baselines}
\end{figure}

Table~\ref{tab:baselines} and Fig.~\ref{fig:baselines} compare the multimodal
final configuration against two text-only baselines under the identical
cross-validation protocol, with all four metrics reported for every row. A
linear SVM on TF-IDF features reaches 76.30 $\pm$ 1.96 points accuracy, and an
XGBoost classifier on TF-IDF plus auxiliary numeric features reaches
78.95 $\pm$ 1.35 points---both roughly 16--19 percentage points below any
CLIP-containing configuration. This confirms that visual features, not caption
text or its derived numeric properties, provide the primary discriminative
signal for this manipulation-detection task. We include the CLIP-only row in
this table deliberately: it is the baseline the proposed configuration must
beat, and at present it does not.

\subsection{Key Observations}

\begin{itemize}
\item CLIP visual features alone match or exceed every text-inclusive
configuration on this dataset---the dataset's manipulations are photographic,
not text-embedded, so visual features dominate.

\item OCR contributes no measurable benefit on BanFakeMM specifically; this is
not generalized as a claim about Bangla misinformation broadly, since
MultiBanFakeDetect \cite{multibanfake} indicates text-in-image remains a
meaningful vector in this language more broadly.

\item The distribution-shift gap (95.40 $\rightarrow$ 86.0 percent accuracy,
95.75 $\rightarrow$ 73.0 percent fake-recall) is substantial even for a
well-performing same-distribution classifier, reinforcing that
same-distribution accuracy alone should not be reported as a proxy for
real-world reliability.

\item Threshold calibration alone appears to recover most of the gap
(86.0 $\rightarrow$ 90.0 percent), with training-data diversification
contributing a further, smaller gain (90.0 $\rightarrow$ 91.5 percent)---but
the first of these is measured on the same set used to select the threshold,
and the second measures domain adaptation rather than generalization.
\end{itemize}

\subsection{Verification Layer Behaviour}

Unlike the classifier, the verification layer is evaluated qualitatively
through targeted case studies covering in-scope claims (relevant fact-check
evidence exists in the local corpus) and out-of-scope claims (none exists),
rather than through a live deployment trial or a quantitative benchmark. These
case studies were not sampled according to a documented procedure and do not
report inter-annotator agreement, so they are illustrative of the system's
operation under both conditions rather than a validated measurement of
verification accuracy or reliability.

A local demonstration application implementing this pipeline was used to
exercise both the direct-image-match and dual-LLM-verification paths during
development, but no live, sustained-monitoring deployment trial---comparable
to a production rollout---was conducted for this paper; this is stated
explicitly rather than implied, and a formal quantitative evaluation of the
verification layer is identified as future work in
Section~\ref{sec:limitations}.

% TODO(reviewer-5): add a worked-examples table here -- 2-3 real claims with
% the retrieved evidence, both models' verdicts, and the consensus outcome.
% The reviewer explicitly asked for "actual examples"; none are currently
% shown anywhere in the paper.

%==============================================================
\section{Discussion}
%==============================================================

\subsection{Visual-versus-Textual Modality Trade-off Analysis}

The ablation outcomes are explained by the nature of BanFakeMM's manipulations
rather than by any inherent weakness in text-based features generally. The
dataset's fake images are produced through pixel-level editing or full-image
AI generation, so the evidence of manipulation is encoded directly in visual
texture, lighting consistency, and generation artifacts---exactly the signal
CLIP's visual encoder is trained to capture. A caption or OCR-derived text
feature can only reflect manipulation indirectly, which is a much weaker and
noisier proxy; this explains both OCR's near-chance standalone performance and
CLIP's dominance once included. This also bounds the finding: on a corpus
where manipulation is carried by embedded text, the ranking could invert.

The absence of a measurable CISF contribution on top of CLIP is similarly
explainable rather than anomalous. CISF's neighbour-weighted fusion mechanism
is designed to sharpen a semantic embedding by borrowing signal from similar
items in the training corpus---a mechanism suited to text embeddings that
benefit from smoothing against semantically related neighbours. One plausible
explanation for CISF's limited effect on CLIP is that CLIP's visual embedding
already encodes a rich, high-dimensional description of the image directly
from pixels, leaving comparatively little redundant noise for neighbour-fusion
to average away---though this is inferred from the observed ablation pattern
rather than independently verified.

Both before and after mitigation, the shifted classifier's errors are
concentrated at the boundary between the training distribution's manipulation
style (localized AI editing) and the evaluation set's manipulation style
(full-image AI generation)---a pattern consistent with, though not proof of, a
decision boundary calibrated tightly around the training distribution's
specific manipulation signature rather than a broader notion of
manipulated-versus-authentic. Whether threshold lowering and diversification
improve the model's underlying visual discrimination ability in a
style-agnostic sense, rather than only shifting the decision boundary, cannot
be determined from these results alone.

\subsection{Addressing the Research Problems}

\begin{itemize}
\item \textbf{Training-data scarcity:} addressed through BanFakeMM's
adversarially constructed training set combined with an explicit, disjoint
distribution-shift evaluation---rather than assuming same-distribution
accuracy generalizes to real deployment.

\item \textbf{Claim-caption divergence:} addressed by supporting
user-supplied claim text as a first-priority input to the verification layer,
with the BLIP caption used only as a documented fallback.

\item \textbf{Evidence-source scarcity:} addressed by building a dedicated
Bangla-language evidence corpus, while explicitly acknowledging its
single-publisher composition as a scoping limitation rather than a claim of
broad coverage.

\item \textbf{Missing verification stage:} addressed by adding a
retrieval-augmented, dual-independent-LLM verification layer with direct
image-to-image evidence matching, evaluated qualitatively and scoped as a
prototype rather than left unaddressed or overstated.

\item \textbf{Same-distribution overstatement:} addressed by reporting the
zero-shot shifted result (86.0 percent accuracy, 73.0 percent fake-recall)
alongside the same-distribution result (95.40 percent) rather than reporting
only the higher, more favourable number---and by labelling the post-adaptation
figure as domain adaptation rather than as recovered generalization.
\end{itemize}

\subsection{Verification Layer Reliability Considerations}

Because the verification layer's central role is evidence-grounded claim
assessment, its reliability considerations differ from the classifier's. Two
design choices bound how much confidence should be placed in its output today.
First, the dual-LLM consensus rule is deliberately conservative: a verdict is
returned only when both models complete and agree, and any disagreement or
failure is routed to human review rather than silently resolved by one model's
output---this avoids overstating confidence, at the cost of leaving more cases
unresolved automatically. Second, the 0.90 image-similarity threshold used for
direct-match evidence was set from an informal empirical separation observed
during development, not a labelled precision/recall study, so it should be
treated as a practical operating point rather than a statistically validated
boundary. Both choices trade automation for caution; a production deployment
intending to reduce the human-review rate would need the quantitative
verification-layer evaluation identified in Section~\ref{sec:limitations}
before relaxing either safeguard.

\subsection{Limitations and Future Directions}
\label{sec:limitations}

Several limitations bound the scope of the present work, in addition to the
two disclosures already stated in Sections~\ref{sec:verify} and
\ref{sec:shift}.

Ground-truth labels throughout this study derive from BanFakeMM's own
deterministic construction process, so the classifier's demonstrated
competence is specific to detecting this class of algorithmic image
manipulation, not to adjudicating whether a news item's overall claim is
true---a task the verification layer, not the classifier, is meant to address.
The distribution-shift evaluation, while genuinely disjoint from training
data, is drawn from a second split of the same underlying benchmark;
generalization to manipulation styles and image sources entirely outside
BanFakeMM, including organically collected real-world Bangladeshi social media
imagery, remains an open question.

The evidence-verification layer's evaluation in this paper is qualitative
rather than quantitative. A formal evaluation protocol---retrieval recall@$k$
against a labelled evidence set, verdict accuracy against ground-truth-labeled
claims---is left to future work; established RAG-evaluation frameworks such as
RAGProbe \cite{ragprobe} and the faithfulness, context-relevance and
answer-relevance metrics summarized by Roychowdhury et al. \cite{ragmetrics}
are a concrete starting point, alongside expanding the evidence corpus beyond
its current single fact-checking source.

The threshold-calibration and data-diversification approach used to narrow the
distribution-shift gap is a data-centric mitigation. Architecture-level
alternatives---CLIPood's fine-tuning-based approach to CLIP out-of-distribution
generalization \cite{clipood}, domain-specific prompt-tuning approaches such as
DPOD \cite{dpod}---were not evaluated here. Threshold-independent robustness
metrics such as AUROC and FPR@95, and calibration metrics such as expected
calibration error and Brier score, were likewise not computed for this
evaluation, and cross-validation folds were not explicitly grouped by source
article or near-duplicate image, which is noted as an unresolved leakage risk
rather than a ruled-out one. Finally, this pipeline addresses static image
manipulation; extending it to video-based misinformation would require a
substantially different technical approach beyond the scope of the present
system.

%==============================================================
\section{Conclusion}
%==============================================================

Practical Bangla misinformation detection demands more than a
same-distribution accuracy figure---it requires an explicit distribution-shift
check, an honest account of which modalities actually contribute, and, where a
classification score alone is insufficient, an evidence-grounded verification
stage.

This research contributes: (1) a multimodal Bangla fake-news image classifier
combining BLIP, SBERT/CISF, and CLIP features, with its final configuration
selected through per-modality ablation rather than default inclusion of every
signal; (2) an empirically justified architectural simplification---OCR
removed after ablation showed it performs near chance level on this dataset;
(3) a distribution-shift evaluation protocol quantifying a real generalization
gap (95.40 $\rightarrow$ 86.0 percent accuracy, 95.75 $\rightarrow$
73.0 percent fake-recall) and, separately, the extent to which supervised
domain adaptation on target-style data recovers it; (4) a retrieval-augmented,
dual-independent-LLM verification prototype, grounded in a Bangla-language
evidence corpus, that reaches an evidence-based verdict either from claim text
or directly from CLIP-based image-to-image evidence matching; and
(5) transparent reporting of the concrete limitations of both components,
including the qualitative-only status of the verification-layer evaluation and
the threshold-selection effect affecting two rows of
Table~\ref{tab:shift}.

The gap between a high same-distribution accuracy figure and a system that
behaves reliably outside its training distribution is real and, on this task,
substantial. By combining empirical modality selection, a genuine
distribution-shift evaluation, and evidence-grounded verification rather than
classifier-score-only reporting, this work outlines one concrete path toward
Bangla multimodal misinformation detection that is validated against, rather
than merely trained on, the conditions it is meant to operate under.

The immediate next steps follow directly from the disclosures above: a sealed
test partition with threshold selection confined to a disjoint validation
split; near-duplicate verification by perceptual hashing and embedding
similarity with group-aware folds; a paired significance and calibration
comparison to decide whether the fusion pathway is retained or the classifier
simplified to CLIP-only; and a labelled claim set with retrieval and verdict
metrics to promote the verification stage from prototype to evaluated
component. Future work should also extend this evaluation to real-world,
organically collected Bangladeshi social media imagery and broaden the
evidence corpus beyond its current single publisher.

%==============================================================
% ACKNOWLEDGMENT -- hidden while \blindsubmissiontrue, because naming a
% funder, institution or colleague defeats anonymous review. It prints
% automatically once you set \blindsubmissionfalse for the camera-ready.
%==============================================================
\ifblindsubmission\else
\section*{Acknowledgment}
The authors thank [NAMES] for their support of this work. [Add funding
agency and grant number here if applicable; delete this section entirely if
neither applies.]
\fi

%==============================================================
\begin{thebibliography}{00}

% TODO(reviewer-6, BLOCKING): the BanFakeMM record below could not be
% independently verified -- supply curator, year and DOI from the actual
% Mendeley Data landing page before submission.
\bibitem{banfakemm} ``BanFakeMM: Bangla multimodal fake news dataset,''
Mendeley Data. [Online]. Available: \url{https://data.mendeley.com/}
\ifdraftnotes\\\emph{\footnotesize [DRAFT NOTE -- REMOVE BEFORE SUBMISSION]
Dataset citation incomplete: no DOI, curator name, or year could be
independently verified from public sources or this project's own repository;
the exact Mendeley record must be supplied before submission.}\fi

\bibitem{banfakenews} M. Z. Hossain, M. A. Rahman, M. S. Islam, and S. Kar,
``BanFakeNews: A dataset for detecting fake news in Bangla,'' in
\emph{Proc. 12th Lang. Resources and Evaluation Conf. (LREC)}, Marseille,
France, 2020, pp. 2862--2871.

\bibitem{multibanfake} F. T. J. Faria \emph{et al.},
``MultiBanFakeDetect: Integrating advanced fusion techniques for multimodal
detection of Bangla fake news in under-resourced contexts,''
\emph{Int. J. Inf. Manage. Data Insights}, 2025.

\bibitem{banfakenews2} H. M. Shibu, S. Datta, M. S. Miah, N. Sami,
M. S. Chowdhury, and M. S. Islam, ``From scarcity to capability: Empowering
fake news detection in low-resource languages with LLMs,'' in
\emph{Proc. Conf. Indo NLP (IndoNLP)}, 2025, arXiv:2501.09604.

\bibitem{mmcfnd} S. Bansal, N. S. Singh, S. S. Dar, and N. Kumar,
``MMCFND: Multimodal multilingual caption-aware fake news detection for
low-resource Indic languages,'' 2024, arXiv:2410.10407.

\bibitem{eann} Y. Wang \emph{et al.}, ``EANN: Event adversarial neural
networks for multi-modal fake news detection,'' in \emph{Proc. 24th ACM
SIGKDD Int. Conf. Knowl. Discovery Data Mining}, 2018, pp. 849--857.

\bibitem{mvae} D. Khattar, J. S. Goud, M. Gupta, and V. Varma,
``MVAE: Multimodal variational autoencoder for fake news detection,'' in
\emph{Proc. World Wide Web Conf. (WWW)}, San Francisco, CA, 2019,
pp. 2915--2921.

\bibitem{spotfake} S. Singhal, R. R. Shah, T. Chakraborty, P. Kumaraguru, and
S. Satoh, ``SpotFake: A multi-modal framework for fake news detection,'' in
\emph{Proc. IEEE 5th Int. Conf. Multimedia Big Data (BigMM)}, Singapore, 2019,
pp. 39--47.

\bibitem{defame} T. Braun, M. Rothermel, M. Rohrbach, and A. Rohrbach,
``DEFAME: Dynamic evidence-based fact-checking with multimodal experts,''
2024, arXiv:2412.10510.

\bibitem{ragar} M. A. Khaliq, P. Chang, M. Ma, B. Pflugfelder, and
F. Mileti\'{c}, ``RAGAR, your falsehood radar: RAG-augmented reasoning for
political fact-checking using multimodal large language models,'' in
\emph{Proc. 7th FEVER Workshop, ACL}, 2024, arXiv:2404.12065.

% TODO(reviewer-6, BLOCKING): confirm the RAMA author list from the arXiv
% record before submission; it could not be independently verified.
\bibitem{rama} ``RAMA: Retrieval-augmented multi-agent framework for
misinformation detection in multimodal fact-checking,'' 2025,
arXiv:2507.09174.
\ifdraftnotes\\\emph{\footnotesize [DRAFT NOTE -- REMOVE BEFORE SUBMISSION]
Author list not independently confirmed at time of writing; verify against the
arXiv record before submission.}\fi

\bibitem{calm} I.-H. Hsu \emph{et al.}, ``CaLM: Contrasting large and small
language models to verify grounded generation,'' in \emph{Findings Assoc.
Comput. Linguistics: ACL 2024}, arXiv:2406.05365.

\bibitem{mmfusionnet} J. He, T. Liu, J. Zhao, and B. Turner,
``MMFusionNet: Context-aware dynamic fusion for multi-modal fake news
detection with large vision-language models,'' 2025, arXiv:2508.05684.

\bibitem{claaf} G. Mu, C. Chen, X. Li, Y. Chen, J. Dai, and J. Li,
``CLAAF: Multimodal fake information detection based on contrastive learning
and adaptive agg-modality fusion,'' \emph{PLOS ONE}, May 2025.

\bibitem{clipood} Y. Shu, X. Guo, J. Wu, X. Wang, J. Wang, and M. Long,
``CLIPood: Generalizing CLIP to out-of-distributions,'' in \emph{Proc. 40th
Int. Conf. Machine Learning (ICML)}, 2023, arXiv:2302.00864.

\bibitem{dpod} A. Bhattacharya, D. Brahma, S. Nagaje, A. Asati, V. Verma, and
S. Biswas, ``Can out-of-domain data help to learn domain-specific prompts for
multimodal misinformation detection?'' in \emph{Proc. IEEE/CVF Winter Conf.
Appl. Comput. Vis. (WACV)}, 2025, arXiv:2311.16496.

\bibitem{ragmetrics} S. Roychowdhury, S. Soman, H. G. Ranjani, N. Gunda,
V. Chhabra, and S. K. Bala, ``Evaluation of RAG metrics for question answering
in the telecom domain,'' in \emph{ICML 2024 Workshop on Foundation Models in
the Wild}, 2024, arXiv:2407.12873.

\bibitem{ragprobe} S. Sivasothy, S. Barnett, S. Kurniawan, Z. Rasool, and
R. Vasa, ``RAGProbe: An automated approach for evaluating RAG applications,''
2024, arXiv:2409.19019.

\end{thebibliography}

%==============================================================
% TRIM CANDIDATES -- when cutting down to the ICCIT 6-page limit,
% remove in roughly this order (least damaging first):
%
%  1. Fig. 2 (fig1_ablation)      -- Table I has the same numbers
%  2. Fig. 3 (fig2_ood_progression) -- Table II has the same numbers
%  3. Fig. 6 (fig5_baselines)     -- Table IV has the same numbers
%  4. Sec. IV-G "Key Observations" -- repeats IV-A / IV-C / V-A
%  5. Sec. V-B "Addressing the Research Problems" -- repeats Sec. I
%  6. Sec. I-A "Research Problems" bullets -> compress to 3 sentences
%  7. Sec. I-B contribution bullets 4 and 5 -> fold into bullet 3
%  8. Sec. II's four "critical limitations" bullets -> one paragraph
%  9. Sec. IV-B "Final Configuration Evaluation" -> merge into IV-A
% 10. Sec. IV-D confusion-matrix prose -> move into Fig. 5 caption
% 11. Sec. V-C "Verification Layer Reliability" -> merge into III-E Scope
% 12. Sec. IV-H "Verification Layer Behaviour" -> merge into III-E Scope
% 13. Sec. V-A paragraph 2 (CISF explanation)
% 14. Table I rows "Caption only" and "Text only"
%
% Keep at all costs: the three disclosure paragraphs (duplicate
% verification, threshold selection, verification scope), the
% domain-adaptation labelling, and the per-class metrics columns --
% those are the reviewer fixes.
%==============================================================

\end{document}
